\subsection{PHÂN LOẠI VÀ PHƯƠNG PHÁP GIẢI TOÁN}
\begin{dang}{Ghép nhóm mẫu số liệu}
\end{dang}

\begin{vd}%[1K3B8-1]
	Bảng thống kê sau cho biết thời gian chạy (phút) của $30$ vận động viên (VĐV) trong một giải chạy Marathon.
	\begin{center}
		\begin{tabular}{|c|c|c|c|c|c|c|c|c|c|c|c|c|}
			\hline
			TG&$129$&$130$&$133$&$134$&$135$&$136$&$138$&$141$&$142$&$143$&$144$&$145$\\
			\hline
			SL&$1$&$2$&$1$&$1$&$1$&$2$&$3$&$3$&$4$&$5$&$2$&$5$\\
			\hline
		\end{tabular}
	\end{center}
	Hãy chuyển mẫu số liệu trên sang mẫu số liệu ghép nhóm gồm sáu nhóm có độ dài bằng nhau và bằng $3$. Xác định giá trị đại diện của mỗi nhóm.\\
	\dongcham{6}
	\loigiai{
		Giá trị nhỏ nhất là $129$, giá trị lớn nhất là $145$ nên khoảng biến thiên là $145-129=16$. Tổng độ dài của sáu nhóm là $18$. Để cho đối xứng, ta chọn đầu mút trái của nhóm đầu tiên là $127{,}5$ và đầu mút phải của nhóm cuối cùng là $145{,}5$ ta được các nhóm là $[127{,}5;130{,}5),\; [130{,5};133{,5}],\ldots , [142{,}5;145{,}5]$. Đếm số giá trị thuộc mỗi nhóm, ta có mẫu số liệu ghép nhóm như sau
		\begin{center}
			\begin{tabular}{|c|c|c|c|c|c|c|}
				\hline
				Thời gian&$[127{,}5;130{,}5)$&$[130{,}5;133{,}5)$&$[133{,}5;136{,}5)$&$[136{,}5;139{,}5)$&$[139{,}5;142{,}5)$&$[142{,}5;145{,}5)$\\
				\hline
				Số VĐV&$3$&$1$&$4$&$3$&$7$&$12$\\
				\hline
			\end{tabular}
		\end{center}
	}
\end{vd}
\begin{vd}
	Cân nặng (kg) của 35 người trưởng thành tại một khu dân cư được cho như sau:
	\begin{center}
		\begin{tabular}{lllllllllllll}
			43 & 51 & 47 & 62 & 48 & 40 & 50 & 62 & 53 & 56 & 40 & 48 &\\
			56 & 53 & 50 & 42 & 55 & 52 & 48 & 46 & 45 & 54 & 52 & 50 & \\
			47 & 44 & 54 & 55 & 60 & 63 & 58 & 55 & 60 & 58 & 53.
		\end{tabular}
	\end{center}
	Chuyển mẫu số liệu trên thành dạng ghép nhóm, các nhóm có độ dài bằng nhau, trong đó có nhóm $[40; 45)$. Xác định giá trị đại diện của mỗi nhóm.\\
	\dongcham{6}
	\loigiai{
		Vì các nhóm có độ dài bằng nhau, trong đó có nhóm $[40; 45)$ nên độ dài mỗi nhóm là $5$. Ta có mẫu số liệu ghép nhóm
		\begin{center}
			\begin{tabular}{|c|c|c|c|c|c|}
				\hline Cân nặng (kg) &$[40;45)$&$[45;50)$&$[50;55)$&$[55;60)$&$[60;65)$\\
				\hline Số người &5 & 7&11 &7 &5 \\
				\hline
			\end{tabular}
		\end{center}
	}
\end{vd}


\begin{dang}{Tính số trung bình của mẫu số liệu ghép nhóm}
\end{dang}

\begin{vd}%[1T5B1-2]%[1T5B1-3]
	Anh Văn ghi lại cự li 30 lần ném lao của mình ở bảng sau (đơn vị: mét):
	\begin{center}
		\begin{tabular}{|c|c|c|c|c|c|c|c|c|c|}
			\hline $72{,}1$ & $72{,}9$ & $70{,}2$ & $70{,}9$ & $72{,}2$ & $71{,}5$ & $72{,}5$ & $69{,}3$ & $72{,}3$ & $69{,}7$ \\
			\hline $72{,}3$ & $71{,}5$ & $71{,}2$ & $69{,}8$ & $72{,}3$ & $71{,}1$ & $69{,}5$ & $72{,}2$ & $71{,}9$ & $73{,}1$ \\
			\hline $71{,}6$ & $71{,}3$ & $72{,}2$ & $71{,}8$ & $70{,}8$ & $72{,}2$ & $72{,}2$ & $72{,}9$ & $72{,}7$ & $70{,}7$ \\
			\hline
		\end{tabular}
	\end{center}
	\begin{enumerate}
		\item Tính cự li trung bình của mỗi lần ném.
		\item Tổng hợp lại kết quả ném của anh Văn vào bảng tần số ghép nhóm theo mẫu sau:
		\begin{center}
			\begin{tabular}{|c|c|c|c|c|c|}
				\hline Cự li $(\mathrm{m})$ &{$[69{,}2; 70)$} &{$[70; 70{,}8)$} &{$[70{,}8; 71{,}6)$} &{$[71{,}6; 72{,}4)$} &{$[72{,}4; 73{,}2)$} \\
				\hline Số lần & $?$ & $?$ & $?$ & $?$ & $?$ \\
				\hline
			\end{tabular}
		\end{center}
		\item Hãy ước lượng cự li trung bình mỗi lần ném từ bảng tần số ghép nhóm trên.
		% \item Khả năng anh Văn ném được khoảng bao nhiêu mét là cao nhất?
	\end{enumerate}
	\dongcham{7}
	\loigiai{
		\begin{enumerate}
			\item Điểm tổng của mỗi đợt gồm 10 lần ném
			\begin{center}
				\begin{tabular}{|c|c|c|c|c|c|c|c|c|c|c|}
					\hline Điểm &Điểm &Điểm &Điểm &Điểm &Điểm &Điểm &Điểm &Điểm &Điểm &Tổng \\
					\hline $72{,}1$ & $72{,}9$ & $70{,}2$ & $70{,}9$ & $72{,}2$ & $71{,}5$ & $72{,}5$ & $69{,}3$ & $72{,}3$ & $69{,}7$ &$713{,}6$\\
					\hline $72{,}3$ & $71{,}5$ & $71{,}2$ & $69{,}8$ & $72{,}3$ & $71{,}1$ & $69{,}5$ & $72{,}2$ & $71{,}9$ & $73{,}1$ &$714{,}9$\\
					\hline $71{,}6$ & $71{,}3$ & $72{,}2$ & $71{,}8$ & $70{,}8$ & $72{,}2$ & $72{,}2$ & $72{,}9$ & $72{,}7$ & $70{,}7$ &$718{,}4$\\
					\hline
				\end{tabular}
			\end{center}
			Cự li trung bình của mỗi lần ném của anh Văn
			\[\overline{x}=\dfrac{713{,}6+714{,}9+718{,}4}{30}\approx71{,}56\ (\mathrm{m}). \]
			\item Bảng tần số ghép nhóm kết quả ném của anh Văn:
			\begin{center}
				\begin{tabular}{|c|c|c|c|c|c|}
					\hline Cự li $(\mathrm{m})$ &{$[69{,}2; 70)$} &{$[70; 70{,}8)$} &{$[70{,}8; 71{,}6)$} &{$[71{,}6; 72{,}4)$} &{$[72{,}4; 73{,}2)$} \\
					\hline Số lần & $4$ & $2$ & $7$ & $12$ & $5$ \\
					\hline
				\end{tabular}
			\end{center}
			\item Bảng tần số ghép nhóm kết quả ném của anh Văn (theo giá trị đại diện):
			\begin{center}
				\begin{tabular}{|c|c|c|c|c|c|}
					\hline Cự li $(\mathrm{m})$ &{$[69{,}2; 70)$} &{$[70; 70{,}8)$} &{$[70{,}8; 71{,}6)$} &{$[71{,}6; 72{,}4)$} &{$[72{,}4; 73{,}2)$} \\
					\hline Giá trị đại diện &$69{,}6$ &$70{,}4$ &$71{,}2$ &$72{,}0$ &$72{,}8$\\
					\hline Số lần & $4$ & $2$ & $7$ & $12$ & $5$ \\
					\hline
				\end{tabular}
			\end{center}
			Cự li trung bình mỗi lần ném của anh Văn qua bảng tần số ghép nhóm
			\[(69{,}6\cdot 4+70{,}4\cdot 2+71{,}2\cdot 7+72\cdot 12+72{,}8\cdot 5):30=71{,}52\ (\mathrm{m}).  \]
			% \item Nhóm chứa mốt của mẫu số liệu trên là nhóm $[71{,}6; 72{,}4)$.\\
			% Do đó $u_m=71{,}6$; $n_{m-1}=7$; $n_m=12$; $n_{m+1}=5$; $u_{m+1}-u_m=72{,}4-71{,}6=0{,}8$.\\
			% Mốt của mẫu số liệu ghép nhóm là
			% \[M_0=71{,}6+\dfrac{12-7}{(12-7)+(12-5)} \cdot 0{,}8=\dfrac{101}{14} \approx 71{,}93. \]
			% Dựa vào kết quả trên thì khả năng anh Văn ném được cao nhất là khoảng $71{,}93$ mét.
		\end{enumerate}
	}
\end{vd}

\begin{vd}%[1K3B9-1] 
	Tìm cân nặng trung bình của học sinh lớp $11D$ cho trong bảng sau:
	\begin{center}
		\begin{tabular}{|c|c|c|c|c|c|c|}
			\hline
			Cân nặng	& $\left[40{,}5;45{,}5 \right)$ & $\left[45{,}5;50{,}5 \right)$ & $\left[50{,}5;55{,}5 \right)$ & $\left[55{,}5;60{,}5 \right)$ & $\left[60{,}5;65{,}5 \right)$ & $\left[65{,}5;70{,}5 \right)$ \\
			\hline
			Số học sinh&$10$	& $7$ & $16$ &$4$ & $2$ & $3$ \\
			\hline
		\end{tabular}
	\end{center}
	\dongcham{4}
	\loigiai{
		Trong mỗi khoảng cân nặng, giá trị đại diện là trung bình cộng của hai giá trị đầu mút nên ta có bảng sau:
		\begin{center}
			\begin{tabular}{|c|c|c|c|c|c|c|}
				\hline
				Cân nặng (kg)	& $43$ & $48$ & $53$ & $58$ & $63$ & $68$ \\
				\hline
				Số học sinh&$10$	& $7$ & $16$ &$4$ & $2$ & $3$ \\
				\hline
			\end{tabular}
		\end{center}	
		Tổng số học sinh là $n=42$. Cân nặng trung bình của học sinh lớp $11D$ là $$\overline{x}=\dfrac{10\cdot 43+7\cdot 48+16\cdot 53+4\cdot 58+2\cdot 63+3\cdot 68}{42}\approx51{,}81\,\mathrm{(kg)}.$$
	}
\end{vd}

\begin{vd}%[1T5B1-2]
	Kết quả khảo sát cân nặng của $25$ quả cam ở mỗi lô hàng $A$ và $B$ được cho ở bảng sau:
	\begin{center}
		\begin{tabular}{|c|c|c|c|c|c|}
			\hline \multicolumn{1}{|c|}{Cân nặng $(\mathrm{g})$} &{$[150; 155)$} &{$[155; 160)$} &{$[160; 165)$} &{$[165; 170)$} &{$[170; 175)$} \\
			\hline Số quả cam ở lô hàng $A$ & 2 & 6 & 12 & 4 & 1 \\
			\hline Số quả cam ở lô hàng $B$ & 1 & 3 & 7 & 10 & 4 \\
			\hline
		\end{tabular}
	\end{center}
	\begin{enumerate}
		\item Hãy ước lượng cân nặng trung bình của mỗi quả cam ở lô hàng $A$ và lô hàng $B$.
		\item Nếu so sánh theo số trung bình thì cam ở lô hàng nào nặng hơn?
	\end{enumerate}
	\dongcham{7}
	\loigiai{
		Ta có bảng thống kê số lượng cam theo giá trị đại diện:
		\begin{center}
			\begin{tabular}{|c|c|c|c|c|c|}
				\hline \multicolumn{1}{|c|}{Cân nặng $(\mathrm{g})$} &{$152{,}5$} &{$157{,}5$} &{$162{,}5$} &{$167{,}5$} &$172{,}5$\\
				\hline Số quả cam ở lô hàng $A$ & 2 & 6 & 12 & 4 & 1 \\
				\hline Số quả cam ở lô hàng $B$ & 1 & 3 & 7 & 10 & 4 \\
				\hline
			\end{tabular}
		\end{center}
		\begin{enumerate}
			\item Cân nặng trung bình của mỗi quả cam ở lô hàng $A$ xấp xỉ bằng
			\[(2\cdot 152{,}5+6\cdot 157{,}5+12\cdot 162{,}5+4\cdot 167{,}5+1\cdot 172{,}5): 25=161{,}7\ (\mathrm{g}). \]
			Cân nặng trung bình của mỗi quả cam ở lô hàng $B$ xấp xỉ bằng
			\[(1\cdot 152{,}5+3\cdot 157{,}5+7\cdot 162{,}5+10\cdot 167{,}5+4\cdot 172{,}5): 25=165{,}1\ (\mathrm{g}). \]
			\item Nếu so sánh theo số trung bình thì cam ở lô hàng $B$ nặng hơn cam ở lô hàng $A$.
		\end{enumerate}
	}
\end{vd}

\begin{dang}{Tính mốt của mẫu số liệu ghép nhóm}
\end{dang}

\begin{vd}%[1K3B9-4]
	Bảng số liệu ghép nhóm sau cho biết chiều cao (cm) của $50$ học sinh lớp $11A$.
	\begin{center}
		\begin{tabular}{|c|c|c|c|c|c|c|}
			\hline
			Khoảng chiều cao (cm)	& $\left[145;150 \right)$ & $\left[150;155 \right)$ & $\left[155;160 \right)$ & $\left[160;165 \right)$&$\left[165;170 \right)$ \\
			\hline
			Số học sinh&$7$	& $14$ & $10$ &$10$ & $9$ \\
			\hline
		\end{tabular}
	\end{center}
	Tính mốt của mẫu số liệu ghép nhóm này. Có thể kết luận gì từ giá trị tính được?	
	\loigiai{
		Tần số lớn nhất là $14$ nên nhóm chứa mốt là nhóm $\left[150;155 \right)$. Ta có $j=2$, $a_2=150$, $m_2=14$, $m_1=7$, $m_3=10$, $h=5$. Do đó $$M_o=150+\dfrac{14-7}{\left(14-7\right)+\left(14-10\right)}\cdot 5\approx 153{,}18.$$	
		Số học sinh có chiều cao khoảng $153{,}18$ là nhiều nhất.
	}
\end{vd}\dongcham{4}
\begin{vd}%[1C5B1-5]
	Kết quả kiểm tra môn Toán của lớp $11D$ như sau
	\[
	\begin{array}{cccccccccccccccccccc}
		5 & 6 & 7 & 5 & 6 & 9 & 10 & 8 & 5 & 5 & 4 & 5 & 4 & 5 & 7 & 4 & 5 & 8 & 9 & 10 \\
		5 & 3 & 5 & 6 & 5 & 7 & 5 & 8 & 4 & 9 & 5 & 6 & 5 & 6 & 8 & 8 & 7 & 9 & 7 & 9
	\end{array}
	\]
	\begin{enumerate}
		\item Lập bảng tần số ghép nhóm của mẫu số liệu trên có bốn nhóm ứng với bốn nửa khoảng $\left[3;5\right)$, $\left[5;7\right)$, $\left[7;9\right)$, $\left[9;11\right)$.
		\item Mốt của bảng số liệu ghép nhóm trên là bao nhiêu (làm tròn kết quả đến hàng phần mười)?
	\end{enumerate}
	\loigiai{
		\immini
		{
			\begin{enumerate}
				\item Bảng bên là bảng tần số ghép nhóm cho kết quả kiểm tra môn Toán của lớp $11D$.
				\item Ta thấy: Nhóm $2$ ứng với nửa khoảng $\left[5;7\right)$ là nhóm có tần số lớn.
				Khi đó, mốt của mẫu số liệu là 
				\[
				M_o = 5 + \dfrac{18- 5}{(18 - 5)+(18- 10)}\cdot 2 \approx 6{,}2.
				\]
			\end{enumerate}
		}
		{
			\begin{tabular}{|c|c|}
				\hline
				\textbf{Nhóm} & \textbf{Tần số}\\ 
				\hline
				$\left[3;5\right)$ & $5$\\
				\hline
				$\left[5;7\right)$ & $18$\\
				\hline
				$\left[7;9\right)$ & $10$\\
				\hline
				$\left[9;11\right)$ & $7$\\
				\hline
				& $n = 40$ \\
				\hline
			\end{tabular}
		}
	}
\end{vd}\dongcham{6}
\subsection{BÀI TẬP TỰ LUYỆN}

\begin{bt}%[1C5B1-5]
	Mẫu số liệu sau ghi lại cân nặng của $30$ bạn học sinh (đơn vị: kilôgam):
	\begin{center}
		\begin{tabular}{llllllllll}
			17 & 40 & 39 & 40{,}5 & 42 & 51 & 41{,}5 & 39 & 41 & 30\\
			40 & 42 & 40{,}5 & 39{,}5 & 41 & 40{,}5 & 37 & 39{,}5 & 40 & 41\\
			38{,}5 & 39{,}5 & 40 & 41 & 39 & 40{,}5 & 40 & 38{,}5 & 39{,}5 & 41{,}5
		\end{tabular}
	\end{center}
	Hãy chuyển mẫu số liệu trên sang dạng ghép nhóm với 8 nhóm có độ dài bằng nhau.
	\loigiai{
		Ta có bảng tần số ghép nhóm của mẫu số liệu trên như sau:
		\begin{center}
			\begin{tabular}{|c|c|c|c|c|c|c|c|c|}
				\hline Tốc độ (km/h) &
				$\left[15;20\right)$ &
				$\left[20;25\right)$ &  
				$\left[25;30\right)$ & 
				$\left[30;35\right)$ & 
				$\left[35;40\right)$ &
				$\left[40;45\right)$ & 
				$\left[45;50\right)$ & 
				$\left[50;55\right)$\\
				\hline Số ôtô & 
				1 &
				0 &
				0 &
				1 &
				10 &
				17 &
				0 &
				1 \\
				\hline
			\end{tabular}
		\end{center}
	}
\end{bt}\cham{10}


\begin{bt}%[1T5B1-2]%[1T5B1-3]
	Người ta đếm số xe ô tô đi qua một trạm thu phí mỗi phút trong khoảng thời gian từ $9$ giờ đến $9$ giờ $30$ phút sáng. Kết quả được ghi lại ở bảng sau:
	\begin{center}
		\begin{tabular}{|c|c|c|c|c|c|c|c|c|c|c|c|c|c|c|}
			\hline $15$ & $16$ & $13$ & $21$ & $17$ & $23$ & $15$ & $21$ & $6$ & $11$ & $12$ & $23$ & $19$ & $25$ & $11$ \\
			\hline $25$ & $7$ & $29$ & $10$ & $28$ & $29$ & $24$ & $6$ & $11$ & $23$ & $11$ & $21$ & $9$ & $27$ & $15$ \\
			\hline
		\end{tabular}
	\end{center}
	\begin{enumerate}
		\item Tính số xe trung bình đi qua trạm thu phí trong mỗi phút.
		\item Tổng hợp lại số liệu trên vào bảng tần số ghép nhóm theo mẫu sau:
		\begin{center}
			\begin{tabular}{|c|c|c|c|c|c|}
				\hline Số xe &{$[6; 10]$} &{$[11; 15]$} &{$[16; 20]$} &{$[21; 25]$} &{$[26; 30]$} \\
				\hline Số lần & $?$ & $?$ & $?$ & $?$ & $?$ \\
				\hline
			\end{tabular}
		\end{center}
		\item Hãy ước lượng trung bình số xe đi qua trạm thu phí trong mỗi phút từ bảng tần số ghép nhóm trên.
	\end{enumerate}
	\loigiai{
		\begin{enumerate}
			\item Bảng tần số
			\begin{center}
				\begin{tabular}{|c|c|c|c|c|c|c|c|c|c|c|c|c|c|c|c|c|c|c|c|}
					\hline Giá trị &$6$ & $7$ & $9$ & $10$ & $11$ & $12$ & $13$ & $15$ & $16$ & $17$ & $19$ & $21$ & $23$ & $24$ & $25$ & $27$ & $28$ & $29$ &\\
					\hline Tần số &$2$ & $1$ & $1$ & $1$ & $4$ & $1$ & $1$ & $3$ & $1$ & $1$ & $1$ & $3$ & $3$ & $1$ & $2$ & $1$ & $1$ & $2$ &$N=30$\\
					\hline
				\end{tabular}
			\end{center}
			Số xe trung bình đi qua trạm thu phí trong mỗi phút là
			\allowdisplaybreaks
			\begin{eqnarray*}
				\overline{x}&=&\dfrac{6\cdot 2+7+9+10+11\cdot 4+12+13+15\cdot 3}{30}\\
				&&+\dfrac{16+17+19+21\cdot 3+23\cdot 3+24+25\cdot 2+27+28+29\cdot 2}{30}\\
				&\approx& 17{,}43\ (\text{xe}).
			\end{eqnarray*}	
			\item Bảng tần số ghép nhóm
			\begin{center}
				\begin{tabular}{|c|c|c|c|c|c|}
					\hline Số xe &{$[6; 10]$} &{$[11; 15]$} &{$[16; 20]$} &{$[21; 25]$} &{$[26; 30]$} \\
					\hline Số lần & $5$ & $9$ & $3$ & $9$ & $4$ \\
					\hline
				\end{tabular}
			\end{center}
			\item Bảng tần số ghép nhóm (theo giá trị đại diện) được hiệu chỉnh lại như sau
			\begin{center}
				\begin{tabular}{|c|c|c|c|c|c|}
					\hline Số xe &{$[5{,}5; 10{,}5)$} &{$[10{,}5; 15{,}5)$} &{$[15{,}5; 20{,}5)$} &{$[20{,}5; 25{,}5)$} &{$[25{,}5; 30{,}5)$} \\
					\hline Giá trị đại diện &{$8$} &{$13$} &{$18$} &{$23$} &{$28$} \\
					\hline Số lần & $5$ & $9$ & $3$ & $9$ & $4$ \\
					\hline
				\end{tabular}
			\end{center}
			Số xe trung bình đi qua trạm qua bảng tần số ghép nhóm là
			\[\overline{x}=\dfrac{8\cdot 5+13\cdot 9+18\cdot 3+23\cdot 9+28\cdot 4}{30}\approx 17{,}67\ (\text{xe}). \]
		\end{enumerate}
	}
\end{bt}\cham{16}

\begin{bt}%[1T5B1-2]%[1T5B1-3]
	Một thư viện thống kê số lượng sách được mượn mỗi ngày trong ba tháng ở bảng sau:
	\begin{center}
		\begin{tabular}{|c|c|c|c|c|c|c|c|}
			\hline Số sách &{$[16; 20]$} &{$[21; 25]$} &{$[26; 30]$} &{$[31; 35]$} &{$[36; 40]$} &{$[41; 45]$} &{$[46; 50]$} \\
			\hline Số ngày & 3 & 6 & 15 & 27 & 22 & 14 & 5 \\
			\hline
		\end{tabular}
	\end{center}
	Hãy ước lượng số trung bình và mốt của mẫu số liệu ghép nhóm trên.
	\loigiai{
		Vì số lượng sách được mượn là số nguyên nên ta hiệu chỉnh bảng tần số ghép nhóm (theo giá trị đại diện) như sau
		\begin{center}
			{\footnotesize \begin{tabular}{|c|c|c|c|c|c|c|c|}
					\hline Số sách &{$[15{,}5; 20{,}5)$} &{$[20{,}5; 25{,}5)$} &{$[25{,}5; 30{,}5)$} &{$[30{,}5; 35{,}5)$} &{$[35{,}5; 40{,}5]$} &{$[40{,}5; 45{,}5)$} &{$[45{,}5; 50{,}5)$} \\
					\hline Giá trị đại diện &{$18$} &{$23$} &{$28$} &{$33$} &{$38$} &{$43$} &{$48$} \\
					\hline Số ngày & 3 & 6 & 15 & 27 & 22 & 14 & 5 \\
					\hline
			\end{tabular}}
		\end{center}
		Trung bình số lượng sách được mượn mỗi ngày trong 3 tháng của thư viện là
		\[\overline{x}=\dfrac{18\cdot 3+23\cdot 6+28\cdot 15+33\cdot 27+38\cdot 22+43\cdot 14+48\cdot 5}{92}\approx 34{,}58. \]
		Nhóm chứa mốt của mẫu số liệu trên là nhóm $[30{,}5; 35{,}5)$.\\
		Do đó mốt của mẫu số liệu ghép nhóm là
		\[M_0=30{,}5+\dfrac{27-15}{(27-15)+(27-22)} \cdot 5\approx 34{,}03. \]
	}
\end{bt}\cham{18}

\begin{bt}%[1T5B1-3]
	Một công ty xây dựng khảo sát khách hàng xem họ có nhu cầu mua nhà ở mức giá nào. Kết quả khảo sát được ghi lại ở bảng sau
	\begin{center}
		\begin{tabular}{|c|c|c|c|c|c|}
			\hline \begin{tabular}{c}
				\textbf{Mức giá} \\
				\textbf{(triệu đồng/$\mathrm{m}^2)$}
			\end{tabular} &{$[10; 14)$} &{$[14; 18)$} &{$[18; 22)$} &{$[22; 26)$} &{$[26; 30)$} \\
			\hline \textbf{Số khách hàng} & 54 & 78 & 120 & 45 & 12 \\
			\hline
		\end{tabular}
	\end{center}
	\begin{enumerate}
		\item Tìm mốt của mẫu số liệu ghép nhóm trên.
		\item Công ty nên xây nhà ở mức giá nào để nhiều người có nhu cầu mua nhất?
	\end{enumerate}
	\loigiai{
		\begin{enumerate}
			\item Nhóm chứa mốt của mẫu số liệu trên là nhóm $[18; 22)$.\\ Do đó mốt của mẫu số liệu ghép nhóm là
			\[M_0=18+\dfrac{120-78}{(120-78)+(120-45)} \cdot 4=\dfrac{758}{39} \approx 19{,}4. \]
			\item Dựa vào kết quả trên ta có thể dự đoán rằng nếu công ty xây nhà ở mức giá $19{,}4$ triệu đồng$/ \mathrm{m}^2$ thì sẽ có nhiều người có nhu cầu mua nhất.
		\end{enumerate}
	}
\end{bt}\cham{15}
%%=====Ví dụ 6
\begin{bt}%[1T5B1-3]
	Số cuộc gọi điện thoại một nguời thực hiện mỗi ngày trong $30$ ngày được lựa chọn ngẫu nhiên được thống kê trong bảng sau:
	\begin{center}
		\begin{tabular}{|c|c|c|c|c|c|}
			\hline Số cuộc gọi &{$[3; 5]$} &{$[6; 8]$} &{$[9; 11]$} &{$[12; 14]$} &{$[15; 17]$} \\
			\hline Số ngày & 5 & 13 & 7 & 3 & 2 \\
			\hline
		\end{tabular}
	\end{center}
	\begin{enumerate}
		\item Tìm mốt của mẫu số liệu ghép nhóm trên.
		\item Hãy dự đoán xem khả năng người đó thực hiện bao nhiêu cuộc gọi mỗi ngày là cao nhất.
	\end{enumerate}
	\loigiai{
		Do số cuộc gọi là số nguyên nên ta hiệu chỉnh lại như sau:
		\begin{center}
			\begin{tabular}{|c|c|c|c|c|c|}
				\hline Số cuộc gọi &{$[2{,}5; 5{,}5)$} &{$[5{,}5; 8{,}5)$} &{$[8{,}5; 11{,}5)$} &{$[11{,}5; 14{,}5)$} &{$[14{,}5; 17{,}5)$} \\
				\hline Số ngày & 5 & 13 & 7 & 3 & 2 \\
				\hline
			\end{tabular}
		\end{center}
		\begin{enumerate}
			\item Nhóm chứa mốt của mẫu số liệu trên là nhóm $[5{,}5; 8{,}5)$.\\
			Do đó mốt của mẫu số liệu ghép nhóm là
			\[M_0=5{,}5+\dfrac{13-5}{(13-5)+(13-7)} \cdot 3=\dfrac{101}{14} \approx 7{,}2. \]
			\item Dựa vào kết quả trên ta có thể dự đoán rằng khả năng người đó thực hiện 7 cuộc gọi mỗi ngày là cao nhất.
		\end{enumerate}
	}
\end{bt}\cham{16}
